\renewcommand{\Anno}{2025}

\renewcommand{\NumeroProblema}{1}

\renewcommand{\Testo}{
Nel gioco del {\it Risiko} in ogni battaglia ci sono un attaccante e un difensore, e l'esito viene deciso attraverso il lancio di dadi. Supponiamo che il difensore lanci un solo dado; il suo punteggio è il numero indicato dal dado lanciato. L'attaccante invece lancia $n$ dadi ($n=1,2,3$ nel gioco reale) ed il suo punteggio è il massimo tra i numeri indicati dagli $n$ dadi lanciati. Vince la battaglia chi ha il punteggio più alto; il vantaggio del difensore è che vince la battaglia anche nel caso di parità dei punteggi.

\begin{enumerate}
\renewcommand{\labelenumi}{(\alph{enumi})}

\item Nel caso $n=1$, calcolare la probabilità $p_1$ che il difensore vinca la battaglia.

\item Nel caso $n$ intero positivo qualsiasi, calcolare la probabilità $p_n$ che il difensore vinca la battaglia.

\item Calcolare il limite
$$\lim_{n\to\infty}p_n.$$

\end{enumerate}
}

\renewcommand{\Risposta}{
% SSSUP-BOZZA-AI: Risposta; autore=ChatGPT; data=2026-11-08; stato=DA-VERIFICARE
% SSSUP-PROVENIENZA: AI: rielaborazione e sviluppo di appunti forniti
\begin{enumerate}
\item[(a)] $p_1=7/12$.
\item[(b)] $\displaystyle p_n=\frac1{6^{n+1}}\sum_{k=1}^{6}k^n$.
\item[(c)] $\displaystyle\lim_{n\to\infty}p_n=1/6$.
\end{enumerate}
Le probabilità si intendono per dadi equi e lanci indipendenti.
}

\renewcommand{\Aiutino}{
% SSSUP-BOZZA-AI: Aiutino; autore=ChatGPT; data=2026-11-08; stato=DA-VERIFICARE
% SSSUP-PROVENIENZA: AI: rielaborazione e sviluppo di appunti forniti
Condiziona sul risultato $k$ del difensore. Per vincere, nessuno degli $n$ dadi dell'attaccante deve superare $k$. Nel limite, separa il termine con $k=6$ dagli altri cinque.
}

\renewcommand{\SoluzioneUfficiale}{}

\renewcommand{\SoluzioniAlternative}{}

\renewcommand{\Osservazioni}{}

\renewcommand{\Commenti}{}

% Aggiunta redazionale: la soluzione preesistente non viene modificata.
\providecommand{\SoluzioneRedazionale}{}
\renewcommand{\SoluzioneRedazionale}{
% SSSUP-BOZZA-AI: SoluzioneRedazionale; autore=ChatGPT; data=2026-11-08; stato=DA-VERIFICARE
% SSSUP-PROVENIENZA: AI: rielaborazione e sviluppo di appunti forniti
Assumiamo dadi equi e lanci indipendenti, come nella bozza fornita. Condizionando sul risultato $k$ del difensore, egli vince se e solo se tutti i dadi dell'attaccante mostrano un valore non superiore a $k$. La probabilità condizionata è $(k/6)^n$.

\paragraph{Punto (a).}
Per un solo dado dell'attaccante,
\[
p_1=\frac16\sum_{k=1}^{6}\frac k6=\frac{21}{36}=\frac7{12}.
\]

\paragraph{Punto (b).}
Per ogni intero positivo $n$,
\[
p_n=\frac16\sum_{k=1}^{6}\left(\frac k6\right)^n
=\frac1{6^{n+1}}\sum_{k=1}^{6}k^n.
\]

\paragraph{Punto (c).}
Il termine con $k=6$ vale sempre $1/6$. Gli altri soddisfano
\[
0\leq p_n-\frac16
\leq\frac56\left(\frac56\right)^n
=\left(\frac56\right)^{n+1}\longrightarrow0.
\]
Segue $p_n\to1/6$.
}
